% Options for packages loaded elsewhere
% Options for packages loaded elsewhere
\PassOptionsToPackage{unicode}{hyperref}
\PassOptionsToPackage{hyphens}{url}
\PassOptionsToPackage{dvipsnames,svgnames,x11names}{xcolor}
%
\documentclass[
  spanish,
  letterpaper,
  DIV=11,
  numbers=noendperiod]{scrreprt}
\usepackage{xcolor}
\usepackage{amsmath,amssymb}
\setcounter{secnumdepth}{5}
\usepackage{iftex}
\ifPDFTeX
  \usepackage[T1]{fontenc}
  \usepackage[utf8]{inputenc}
  \usepackage{textcomp} % provide euro and other symbols
\else % if luatex or xetex
  \usepackage{unicode-math} % this also loads fontspec
  \defaultfontfeatures{Scale=MatchLowercase}
  \defaultfontfeatures[\rmfamily]{Ligatures=TeX,Scale=1}
\fi
\usepackage{lmodern}
\ifPDFTeX\else
  % xetex/luatex font selection
\fi
% Use upquote if available, for straight quotes in verbatim environments
\IfFileExists{upquote.sty}{\usepackage{upquote}}{}
\IfFileExists{microtype.sty}{% use microtype if available
  \usepackage[]{microtype}
  \UseMicrotypeSet[protrusion]{basicmath} % disable protrusion for tt fonts
}{}
\makeatletter
\@ifundefined{KOMAClassName}{% if non-KOMA class
  \IfFileExists{parskip.sty}{%
    \usepackage{parskip}
  }{% else
    \setlength{\parindent}{0pt}
    \setlength{\parskip}{6pt plus 2pt minus 1pt}}
}{% if KOMA class
  \KOMAoptions{parskip=half}}
\makeatother
% Make \paragraph and \subparagraph free-standing
\makeatletter
\ifx\paragraph\undefined\else
  \let\oldparagraph\paragraph
  \renewcommand{\paragraph}{
    \@ifstar
      \xxxParagraphStar
      \xxxParagraphNoStar
  }
  \newcommand{\xxxParagraphStar}[1]{\oldparagraph*{#1}\mbox{}}
  \newcommand{\xxxParagraphNoStar}[1]{\oldparagraph{#1}\mbox{}}
\fi
\ifx\subparagraph\undefined\else
  \let\oldsubparagraph\subparagraph
  \renewcommand{\subparagraph}{
    \@ifstar
      \xxxSubParagraphStar
      \xxxSubParagraphNoStar
  }
  \newcommand{\xxxSubParagraphStar}[1]{\oldsubparagraph*{#1}\mbox{}}
  \newcommand{\xxxSubParagraphNoStar}[1]{\oldsubparagraph{#1}\mbox{}}
\fi
\makeatother


\usepackage{longtable,booktabs,array}
\newcounter{none} % for unnumbered tables
\usepackage{calc} % for calculating minipage widths
% Correct order of tables after \paragraph or \subparagraph
\usepackage{etoolbox}
\makeatletter
\patchcmd\longtable{\par}{\if@noskipsec\mbox{}\fi\par}{}{}
\makeatother
% Allow footnotes in longtable head/foot
\IfFileExists{footnotehyper.sty}{\usepackage{footnotehyper}}{\usepackage{footnote}}
\makesavenoteenv{longtable}
\usepackage{graphicx}
\makeatletter
\newsavebox\pandoc@box
\newcommand*\pandocbounded[1]{% scales image to fit in text height/width
  \sbox\pandoc@box{#1}%
  \Gscale@div\@tempa{\textheight}{\dimexpr\ht\pandoc@box+\dp\pandoc@box\relax}%
  \Gscale@div\@tempb{\linewidth}{\wd\pandoc@box}%
  \ifdim\@tempb\p@<\@tempa\p@\let\@tempa\@tempb\fi% select the smaller of both
  \ifdim\@tempa\p@<\p@\scalebox{\@tempa}{\usebox\pandoc@box}%
  \else\usebox{\pandoc@box}%
  \fi%
}
% Set default figure placement to htbp
\def\fps@figure{htbp}
\makeatother


% definitions for citeproc citations
\NewDocumentCommand\citeproctext{}{}
\NewDocumentCommand\citeproc{mm}{%
  \begingroup\def\citeproctext{#2}\cite{#1}\endgroup}
\makeatletter
 % allow citations to break across lines
 \let\@cite@ofmt\@firstofone
 % avoid brackets around text for \cite:
 \def\@biblabel#1{}
 \def\@cite#1#2{{#1\if@tempswa , #2\fi}}
\makeatother
\newlength{\cslhangindent}
\setlength{\cslhangindent}{1.5em}
\newlength{\csllabelwidth}
\setlength{\csllabelwidth}{3em}
\newenvironment{CSLReferences}[2] % #1 hanging-indent, #2 entry-spacing
 {\begin{list}{}{%
  \setlength{\itemindent}{0pt}
  \setlength{\leftmargin}{0pt}
  \setlength{\parsep}{0pt}
  % turn on hanging indent if param 1 is 1
  \ifodd #1
   \setlength{\leftmargin}{\cslhangindent}
   \setlength{\itemindent}{-1\cslhangindent}
  \fi
  % set entry spacing
  \setlength{\itemsep}{#2\baselineskip}}}
 {\end{list}}
\usepackage{calc}
\newcommand{\CSLBlock}[1]{\hfill\break\parbox[t]{\linewidth}{\strut\ignorespaces#1\strut}}
\newcommand{\CSLLeftMargin}[1]{\parbox[t]{\csllabelwidth}{\strut#1\strut}}
\newcommand{\CSLRightInline}[1]{\parbox[t]{\linewidth - \csllabelwidth}{\strut#1\strut}}
\newcommand{\CSLIndent}[1]{\hspace{\cslhangindent}#1}

\ifLuaTeX
\usepackage[bidi=basic,shorthands=off]{babel}
\else
\usepackage[bidi=default,shorthands=off]{babel}
\fi
\ifLuaTeX
  \usepackage{selnolig} % disable illegal ligatures
\fi


\setlength{\emergencystretch}{3em} % prevent overfull lines

\providecommand{\tightlist}{%
  \setlength{\itemsep}{0pt}\setlength{\parskip}{0pt}}



 


\usepackage[papersize={210mm,297mm},lmargin=2.5cm,rmargin=2.5cm,top=3.5cm,bottom=3.5cm]{geometry}

\renewcommand{\baselinestretch}{1.25}
\setlength{\parskip}{1mm}
\usepackage{scrextend}
\usepackage{comment}
\usepackage{caption}
\usepackage{subcaption}
\usepackage{amsmath,amsfonts,amssymb,amsthm,amscd}
\usepackage{pdfsync}
\usepackage{latexsym,xcolor,graphicx}
\usepackage[nottoc]{tocbibind}
\usepackage[utf8]{inputenc}					
\usepackage[T1]{fontenc}
\usepackage[spanish,es-lcroman,es-noquoting]{babel}
\usepackage{csquotes}
\usepackage{tikz}
\usepackage{tikz-cd}
\usepackage{appendix}
\usepackage[hidelinks]{hyperref}
\usepackage{enumitem}\setlist{nosep,after=\vspace{1mm},before=\vspace{0mm}}
\usepackage{graphicx,epstopdf}
\usepackage{tcolorbox}
\usepackage{amssymb,cancel}
\usepackage{amsfonts,amsmath,color}
\usepackage{amsthm,mathrsfs}
\usepackage{enumitem,multicol,bbding}
\usepackage[mathscr]{euscript}
\usepackage{fancyhdr}
\usepackage{etoolbox}
\usepackage{csquotes}
\usepackage{etoolbox}
\usepackage{scrextend}

\usepackage{etoolbox}
\newcommand{\zerodisplayskips}{%
  \setlength{\abovedisplayskip}{2mm}%
  \setlength{\belowdisplayskip}{2mm}%
  \setlength{\abovedisplayshortskip}{1mm}%
  \setlength{\belowdisplayshortskip}{1mm}}
\appto{\normalsize}{\zerodisplayskips}
\appto{\small}{\zerodisplayskips}
\appto{\footnotesize}{\zerodisplayskips}

\usepackage[style=ieee]{biblatex}
%\usepackage{biblatex}
\addbibresource{referencias.bib}

\usepackage{titlesec}
\titleformat{\chapter}[display]
{\bfseries\huge}
{\filleft\huge\chaptertitlename~\thechapter}
{0.5cm}
{\titlerule\filright}
[\titlerule]
\titlespacing*{\chapter}{0pt}{-1.5cm}{2cm}
\titlespacing*{\section}{0pt}{0.25cm}{0.25cm}

\patchcmd{\thebibliography}{\section*{\refname}}{}{}{}

% Cabeceira
\usepackage{fancyhdr}
\pagestyle{fancy}
\fancyhf{}
\renewcommand{\chaptermark}[1]{\markboth{\textbf{\thechapter. #1}}{}}% Formato para o capítulo: N. Nome
\renewcommand{\sectionmark}[1]{\markright{\textbf{\thesection. #1}}}% Formato para a sección: N.M. Nome
\fancyhead[LO]{\rightmark}% Nas páxinas impares, parte esquerda do encabezado, aparecerá o nome do capítulo
\fancyhead[RE]{\leftmark}% Nas páxinas pares, parte dereita do encabezado, aparecerá o nombre da sección
\fancyhead[RO,LE]{\textbf{\thepage}} % Números de páxina nas esquinas dos encabezados

\renewcommand{\quote}{\list{}{\rightmargin=\leftmargin\topsep=0pt}\item\relax}

\newcommand{\portada}[3]{
\thispagestyle{empty}
\definecolor{azulUSC}{RGB}{13,38,120}
\includegraphics[width=13cm]{logomath}
\vspace*{2cm}
\centerline{{\Large\bf  Trabajo Fin de Grado}}
\vspace{2.6cm}
{\center{\Huge\bfseries\color{azulUSC} #1}\par}		
\vspace{1cm}
{\Large
\centerline{#2}		
\vspace{6.4cm}
\centerline{\sf #3}        
\vspace{1.25cm}
\centerline{\sf UNIVERSIDADE DE SANTIAGO DE COMPOSTELA}
}
\enlargethispage{1cm}
\clearpage
}

\newcommand{\primeira}[3]{
\thispagestyle{empty}
\enlargethispage{1cm}
\begin{center}\Large
{\sf GRADO DE MATEM\'ATICAS}\\
\bigskip
{\bfseries Trabajo Fin de Grado}\\
\vspace{4cm}
{\bfseries\Huge #1}\\		
\vspace{1.2cm}
{#2}		\par		     
\end{center}
\vspace{6.5cm}
\begin{center}
\Large
{#3}	\par			    
\vspace{1cm}
{\sf UNIVERSIDADE DE SANTIAGO DE COMPOSTELA}
\end{center}}

\newtheorem{theorem}{Teorema}[chapter]
\newtheorem{proposition}[theorem]{Proposición}
\newtheorem{lemma}[theorem]{Lema}
\newtheorem{corollary}[theorem]{Corolario}
\newtheorem{algorithm}[theorem]{Algoritmo}
\newtheorem{axiom}[theorem]{Axioma}
\newtheorem{case}[theorem]{Caso}
\newtheorem{conclusion}[theorem]{Conclusión}
\newtheorem{condition}[theorem]{Condición}
\newtheorem{conjetura}[theorem]{Conjetura}				
\newtheorem{criterion}[theorem]{Criterio}
\newtheorem{ejercicio}[theorem]{Ejercicio}	
			
\theoremstyle{definition}
\newtheorem{definition}[theorem]{Definición}
\newtheorem{ejemplo}[theorem]{Ejemplo}				
\newtheorem{agradecimientos}[theorem]{Agradecimientos}	

\theoremstyle{remark}
\newtheorem{remark}[theorem]{Observación}
\newtheorem{notation}[theorem]{Notación}
\newtheorem{problem}[theorem]{Problema}
\newtheorem{solution}[theorem]{Solución}

\newtheorem*{remark*}{Observación}
\newtheorem*{remarkimp*}{Observación importante}


\addto\captionsspanish{
	\renewcommand{\contentsname}{\'Indice} 
	\renewcommand\appendixname{Anexo}
	\renewcommand\appendixpagename{Anexos}
}


%PAQUETES E COMANDOS EXTRA
\theoremstyle{definition}
\newtheoremstyle{exampstyle}
{2mm} % Space above
{-2mm} % Space below
{} % Body font
{} % Indent amount
{\bfseries} % Theorem head font
{.} % Punctuation after theorem head
{.5em} % Space after theorem head
{} % Theorem head spec (can be left empty, meaning `normal')
\theoremstyle{exampstyle}
\newtheorem{innercustomgeneric}{\customgenericname}
\providecommand{\customgenericname}{}
\newcommand{\newcustomtheorem}[2]{%
  \newenvironment{#1}[1]
  {%
   \renewcommand\customgenericname{#2}%
   \renewcommand\theinnercustomgeneric{##1}%
   \innercustomgeneric
  }
  {\endinnercustomgeneric}
}
\newcustomtheorem{cdefinition}{Definición}

\AfterEndEnvironment{cdefinition}{\noindent\ignorespaces}

\def\N{\ensuremath{\mathbb{N}}} %-- la N de los naturales
\def\R{\ensuremath{\mathbb{R}}} %-- La R de los reales.
\def\C{\ensuremath{\mathbb{C}}} %-- La C de los complejos.
\def\Q{\ensuremath{\mathbb{Q}}}
\def\K{\ensuremath{\mathbb{K}}}

\newcommand{\vertiii}[1]{{\left\vert\kern-0.25ex\left\vert\kern-0.25ex\left\vert #1 
    \right\vert\kern-0.25ex\right\vert\kern-0.25ex\right\vert}}
\KOMAoption{captions}{tableheading}
\makeatletter
\@ifpackageloaded{float}{}{\usepackage{float}}
\floatstyle{plain}
\@ifundefined{c@chapter}{\newfloat{remark}{h}{lorem}}{\newfloat{remark}{h}{lorem}[chapter]}
\floatname{remark}{Observación}
\newcommand*\listofremarks{\listof{remark}{List of Observacións}}
\makeatother
\makeatletter
\@ifpackageloaded{bookmark}{}{\usepackage{bookmark}}
\makeatother
\makeatletter
\@ifpackageloaded{caption}{}{\usepackage{caption}}
\AtBeginDocument{%
\ifdefined\contentsname
  \renewcommand*\contentsname{Tabla de contenidos}
\else
  \newcommand\contentsname{Tabla de contenidos}
\fi
\ifdefined\listfigurename
  \renewcommand*\listfigurename{Listado de Figuras}
\else
  \newcommand\listfigurename{Listado de Figuras}
\fi
\ifdefined\listtablename
  \renewcommand*\listtablename{Listado de Tablas}
\else
  \newcommand\listtablename{Listado de Tablas}
\fi
\ifdefined\figurename
  \renewcommand*\figurename{Figura}
\else
  \newcommand\figurename{Figura}
\fi
\ifdefined\tablename
  \renewcommand*\tablename{Tabla}
\else
  \newcommand\tablename{Tabla}
\fi
}
\@ifpackageloaded{float}{}{\usepackage{float}}
\floatstyle{ruled}
\@ifundefined{c@chapter}{\newfloat{codelisting}{h}{lop}}{\newfloat{codelisting}{h}{lop}[chapter]}
\floatname{codelisting}{Listado}
\newcommand*\listoflistings{\listof{codelisting}{Listado de Listados}}
\usepackage{amsthm}
\theoremstyle{plain}
\newtheorem{proposition}{Proposición}[chapter]
\theoremstyle{definition}
\newtheorem{definition}{Definición}[chapter]
\theoremstyle{plain}
\newtheorem{theorem}{Teorema}[chapter]
\theoremstyle{remark}
\AtBeginDocument{\renewcommand*{\proofname}{Demostración}}
\newtheorem*{remark}{Observación}
\newtheorem*{solution}{Solución}
\newtheorem{refremark}{Observación}[chapter]
\newtheorem{refsolution}{Solución}[chapter]
\makeatother
\makeatletter
\makeatother
\makeatletter
\@ifpackageloaded{caption}{}{\usepackage{caption}}
\@ifpackageloaded{subcaption}{}{\usepackage{subcaption}}
\makeatother
\usepackage{bookmark}
\IfFileExists{xurl.sty}{\usepackage{xurl}}{} % add URL line breaks if available
\urlstyle{same}
\hypersetup{
  pdftitle={Bases de Schauder en espacios de Banach},
  pdfauthor={Iker Rodriguez Rodriguez},
  pdflang={es},
  colorlinks=true,
  linkcolor={blue},
  filecolor={Maroon},
  citecolor={Blue},
  urlcolor={Blue},
  pdfcreator={LaTeX via pandoc}}


\title{Bases de Schauder en espacios de Banach}
\author{Iker Rodriguez Rodriguez}
\date{2026-06-01}
\begin{document}
\maketitle

\renewcommand*\contentsname{Tabla de contenidos}
{
\setcounter{tocdepth}{2}
\tableofcontents
}

\bookmarksetup{startatroot}

\chapter*{Prefacio}\label{prefacio}
\addcontentsline{toc}{chapter}{Prefacio}

\markboth{Prefacio}{Prefacio}

En lo que sigue se expondrá una versión revisada del Trabajo Fin de
Grado presentado en Julio de 2025 en la Facultad de Matemáticas.

\bookmarksetup{startatroot}

\chapter*{Resumen}\label{resumen}
\addcontentsline{toc}{chapter}{Resumen}

\markboth{Resumen}{Resumen}

Estudiaremos los conceptos básicos de las bases de Schauder,
circunscritas en el ámbito del Análisis Funcional. Veremos algunas de
sus propiedades, tipos y utilidades en la caracterización de atributos
de los espacios, centrándonos sobre todo en la reflexividad y en la
estructura interna de los espacios, dando especial importancia a los
espacios de sucesiones como \(c_0\) y los \(\ell_p\), destacando en
especial a \(\ell_1\) entre estos últimos. También daremos algunas notas
relacionadas con la existencia y unicidad de bases de Schauder e
intentaremos mostrar como estas bases son una extensión natural de las
bases de Hamel en espacios finitos y una generalización de las bases de
Hilbert a espacios de Banach más genéricos.

\vspace*{3.5cm}

\bookmarksetup{startatroot}

\chapter*{Abstract}\label{abstract}
\addcontentsline{toc}{chapter}{Abstract}

\markboth{Abstract}{Abstract}

We will study the basic concepts of Schauder basis, mainly related to
the Functional Analysis. We will see some of their properties, types and
utilities in characterization of spaces attributes, focusing on
reflexivity and the internal structure of spaces, giving special treat
to the spaces of sequences like \(c_0\) and \(\ell_p\), laying emphasis
on \(\ell_1\) among latter. Furthermore, we will give some notes related
to existence and uniqueness of Schauder basis and try to show how this
basis are a natural extension of Hamel basis of finite spaces as well as
a generalization of Hilbert basis to more generic Banach spaces.

\clearpage

\thispagestyle{empty}

\bookmarksetup{startatroot}

\chapter*{Introducción}\label{introducciuxf3n}
\addcontentsline{toc}{chapter}{Introducción}

\markboth{Introducción}{Introducción}

\newline

Desde el inicio de la matemática moderna, con Descartes, hemos tenido el
objetivo de asignar a los objetos de nuestro estudio una secuencia de
cifras con las que poder simplificar dichos objetos y embeberlos en el
mundo de los números, tratándolos así con todo el poder teórico de las
matemáticas. Dio inicio esta tentativa con las primeras coordenadas
\textit{cartesianas} ---así designadas en honor a su promotor--- que
movilizaron el motor de la producción matemática haciéndolo más
fructífero en los últimos \(500\) años que en los previos \(3\,300\) que
hay desde las primeras ternas pitagóricas mesopotámicas conocidas.

Es con esa filosofía en mente que surge la idea de realizar tal
tentativa no ya con objetos físicos o con coordenadas finitas, sino con
el poder del infinito. El objetivo es ahora el estudio de objetos mucho
más abstractos: las funciones ---lo cual trajo consigo el problema de
las inifitas dimensiones---.

Al principio, el estudio se reducía a las funciones de
\(\mathcal{C}^\infty\) o infinitamente diferenciables ---las conocidas
como series de Taylor, 1715, que fallaron en generar todo el espacio y
se reducían a las funciones analíticas---. Con el tiempo, la idea fue
madurando y se quiso ampliar el domino de validez a todas las funciones
continuas, aún cuando no se tenía muy claro qué era una función continua
todavía ---de aquí surgieron las series de Fourier, conocidas ya a
principios del siglo XIX y que resultarían en una base no para las
funciones continuas con la norma de la convergencia uniforme, sino para
el espacio \(L^2\) con la norma \emph{euclidiana}
\(\|\!\cdot\!\|_2\)---.

Con el final del siglo XIX y la llegada del XX comenzaron a darse los
mayores avances en este campo. Hilbert, Schmidt y Riesz---entre muchos
otros--- comenzaron el estudio del espacio de las funciones de cuadrado
integrable y consiguieron extraer ciertas nociones geométricas muy
similares a las de los espacios euclidianos de dimensión finita, como
por ejemplo la ortogonalidad. El concepto creció y se estableció la
abstracción de los espacios de Hilbert ---término acuñado por von
Neumann---, cuyo principal cometido y éxito es generalizar los espacios
euclidianos a dimensión infinita respetando la existencia de sistemas
ortonormales que generan la totalidad del espacio. Surgió así el estudio
sistemático de la noción de base en dimensión infinita.

De forma paralela, Haar construyó un sistema ---el conocido como sistema
de Haar--- de funciones que sirven de base no sólo para el espacio
\(L^2\), que ya tenía una base dada por las series de Fourier que además
funcionaba por ser este un espacio de Hilbert, sino de todos los
espacios \(L^p\), con \(1\leq p<\infty\) que no son Hilbertianos. Poco
después, Georg Faber demostró la existencia de una base para el espacio
de las funciones continuas con la norma uniforme.

Todo este impulso histórico se cristalizó en la tesis de \(1923\) de un
joven matemático polaco, Juliusz Schauder. Él dio rigor y formalismo a
la idea de una generalización de todo lo anterior introduciendo el
concepto hoy en día conocido, en su honor, como
\emph{bases de Schauder}: un avance más en el entendimiento de la
estructura de los espacios de funciones inspirado en el objetivo de
Descartes de simplificar el estudio de cualquier ente a un sistema de
coordenadas coherente del que se puedan extraer conclusiones.

Es por querer dar valor a esta idea tan matemática que en el presente
Trabajo Fin de Grado queremos dar un repaso a algunos de los resultados
que las bases de Schauder han ayudado a probar y comprender.

En el primer capítulo nos centraremos en la definición de lo que es una
base de Schauder y caracterizaremos su existencia; también incluiremos
algunos comentarios sobre el problema de la existencia de base para un
espacio de Banach.

En el Capítulo 2 estudiaremos la relación entre las bases de Schauder y
el espacio dual, demostrando el primer gran resultado del trabajo: la
caracterización de aquellos espacios de Banach que son reflexivos y
admiten base de Schauder.

En el Capítulo 3 estudiaremos las conocidas como sucesiones básicas, que
no son más que bases de Schauder del espacio ---quizás menor que el
total--- que generan. También hablaremos sobre la noción equivalencia
para bases de Schauder.

En el Capítulo 4 estudiaremos aquellas bases de Schauder que son
incondicionales; es decir, es indiferente el orden que consideremos a la
hora de estudiar la convergencia de la serie que define a cada vector
del espacio. Caracterizaremos los espacios de Banach reflexivos con base
incondicional.

En el Capítulo 5 daremos por finalizada nuestra tentativa empleando la
teoría de las bases de Schauder desarrollada a lo largo del trabajo para
estudiar la estructura interna de los espacios de sucesiones \(\ell_p\)
con \(1\leq p<\infty\), con el objetivo de dar a entender como de
distintos son entre sí.

A mayores, queremos dejar patente que las bases de Schauder no son en
absoluto una construcción artificiosa, sino que generalizan muy
naturalmente a las bases de Hamel de los espacios de dimensión finita.
Es por ello que iremos comparando sus propiedades en algunos capítulos
para dejar claro el origen algebraico que tiene detrás esta idea y que
tan bien resume la filosofía del Análisis Funcional de fusionar
conceptos algebraicos con nociones topológicas para el estudio de las
funciones.

\begin{flushright}
    Santiago de Compostela, julio de 2025.
\end{flushright}

\bookmarksetup{startatroot}

\chapter{Bases de Schauder}\label{bases-de-schauder}

En este capítulo presentaremos algunos conceptos básicos: introduciremos
la noción de base de Schauder para un espacio vectorial normado,
caracterizaremos su construcción y existencia y más adelante, ya en la
segunda sección del capítulo, veremos que la hipótesis de completitud es
fundamental para obtener los resultados más importantes de este Trabajo
Fin de Grado.

También hablaremos, aunque de manera más superficial, del \emph{problema
de existencia de base de Schauder} para espacios de Banach separables.

\section{Bases de Schauder en espacios vectoriales
normados.}\label{bases-de-schauder-en-espacios-vectoriales-normados.}

Es bien conocido que todo espacio vectorial tiene base de Hamel. Es
decir, en todo espacio vectorial podemos encontrar un conjunto de
vectores linealmente independientes cuyas combinaciones lineales
(¡finitas!) generan el espacio. Además, todas las bases de Hamel de un
mismo espacio vectorial tienen algo en común: su cardinal. Por tanto, es
posible definir la \textbf{\emph{dimensión}} de cualquier espacio
vectorial.

No obstante, en el caso de dimensión infinita, la exigencia de finitud
de las combinaciones lineales provoca que las bases de Hamel pierdan
gran parte de su utilidad. Más exactamente (véase {[}1{]}, pág. 48):

\begin{theorem}[]\protect\hypertarget{thm-first}{}\label{thm-first}

Ningún espacio vectorial de dimensión infinita admite base de Hamel
numerable.

\end{theorem}

Evidentemente, la no numerabilidad de una base es un gran inconveniente,
pues reduce enormemente la capacidad de trabajar con éstas: con una base
nosotros queremos reducir el número de elementos que necesitamos conocer
de un espacio para poder trabajar con él, pero si dicha base es no
numerable entonces es necesario conocer casi tantos elementos como tiene
el espacio, lo cual hace redundante el sentido de una base.

Para superar la problemática anteriormente mencionada podemos buscar
ayuda en la topología y emplear el concepto de convergencia para
generalizar las combinaciones lineales finitas a sumas infinitas, es
decir, a series.

Sea \((X,\|\!\cdot\!\|)\) un espacio vectorial normado.

\begin{definition}[Base de
Schauder]\protect\hypertarget{def-bs}{}\label{def-bs}

~

\hspace*{1mm}

Se dice que una sucesión \((e_i)_i\subset X\) es una \textbf{\emph{base
de Schauder de \(X\)}} si para cada vector \(x\in X\) existe una única
sucesión de escalares \((a_i)_i\) de modo que\\
\[
x=\sum_{i=1}^{\infty}{a_ie_i},
\]

es decir, cumpliendo que\\
\[
\lim_{n\rightarrow\infty}\big\|x-\sum_{i=1}^{n}{a_ie_i}\big\|=0.
\]

\end{definition}

La anterior definición ya nos deja entrever esa similitud de las bases
de Schauder con las de Hamel: un conjunto de vectores---linealmente
independientes--- que generan el espacio y nos permiten asignar de forma
única unos coeficientes a todos los elementos de dicho espacio. Hasta
aquí tenemos la filosofía algebraica del Análisis Funcional.

El lector podrá notar ahora los beneficios de introducir una perspectiva
analítica, o más bien topológica, al poder reducir el tamaño de las
bases. Donde antes teníamos bases no numerables para espacios de
dimensión infinita, ahora tenemos una sucesión numerable que genera todo
el espacio manteniendo la unicidad de los coeficientes (que tienen
entonces derecho a ser denominados \textbf{\emph{coordenadas}}). Parece
que vamos por buen camino en el cumplimiento de nuestro objetivo.

Veamos ahora una caracterización de estas bases que nos será muy útil
tanto para demostrar que una sucesión es, en efecto, base de Schauder
como para construirlas.

\begin{definition}[Proyecciones canónicas de una base de
Schauder]\protect\hypertarget{def-ProyeccionesCanonicas}{}\label{def-ProyeccionesCanonicas}

~

\hspace{1mm}

Si \((e_i)_i\) es una base de Schauder de \(X\), la sucesión de
\emph{proyecciones canónicas} asociada a \((e_i)_i\) es la sucesión de
operadores \((P_n)_n\subset\mathcal{B}(X)\) dada, para cada
\(n\in\mathbb{N}\), por
\begin{equation}\protect\phantomsection\label{eq-ProyCanon}{
    P_n\colon x=\sum_{i=1}^\infty a_ie_i\in X\longmapsto P_nx=\sum_{i=1}^n a_ie_i\in X.
}\end{equation}

\end{definition}

Sea \((X,\|\!\cdot\!\|)\) un espacio vectorial normado con base de
Schauder \((e_i)_i\).

\begin{proposition}[Caracterización de las proyecciones
canónicas]\protect\hypertarget{prp-proyeciones}{}\label{prp-proyeciones}

~

\hspace*{1mm}

La sucesión de proyecciones canónicas asociada a \((e_i)_i\) satisface
las siguientes propiedades:

\begin{enumerate}
\def\labelenumi{(\alph{enumi})}
\item
  para cada \(n\in\mathbb{N}\), \(P_n(X)\) es un espacio vectorial de
  dimensión \(n\);
\item
  para cualesquiera \(n,\,m\in\mathbb{N}\),
  \(P_nP_m=P_mP_n=P_{\min\{n,m\}}\);
\item
  para cada \(x\in X\), \(\|P_nx-x\|\to 0\) cuando \(n\to\infty\).
\end{enumerate}

Además, dada una sucesión de operadores \((P_n)_n\subset\mathcal{B}(X)\)
satisfaciendo \((a)\), \((b)\) y \((c)\), existe una base de Schauder de
\(X\) que tiene a \((P_n)_n\) como sucesión de proyecciones canónicas
asociada.

\end{proposition}

Antes de ver la demostración de la Proposición
Proposición~\ref{prp-proyeciones} conviene recordar que se dice que un
operador \(P\in\mathcal{B}(X)\) es una \textbf{\emph{proyección}}
(lineal y continua) si \(P^2\equiv PP=P\). En tal caso, \(I-P\) también
es una proyección. Además, \(P(X)={\rm ker}\,(I-P)\), luego \(P\)
restringido a \(P(X)\) es la identidad. Análogamente,
\({\rm ker}\,P=(I-P)(X)\). Es evidente que si \(P\) es una proyección no
trivial (\(P\neq 0\) y \(P\neq I\)), entonces \(\|P\|\geq 1\) y
\(\|I-P\|\geq 1\).

\bookmarksetup{startatroot}

\chapter*{References}\label{references}
\addcontentsline{toc}{chapter}{References}

\markboth{References}{References}

\protect\phantomsection\label{refs}
\begin{CSLReferences}{0}{0}
\bibitem[\citeproctext]{ref-BST}
\CSLLeftMargin{{[}1{]} }%
\CSLRightInline{Marián. Fabian, Petr. Habala, Petr. Hájek, Vicente.
Montesinos, y Václav. Zizler, \emph{Banach Space Theory: The Basis for
Linear and Nonlinear Analysis}. Springer, 2011.}

\end{CSLReferences}




\end{document}
